\appendix
\section*{Appendix H: Holonomy and Projection Geometry in the Altitude System}

\subsection*{H.0 Overview}

This appendix presents the geometric-topological formalization of the
altitude architecture used in the Constellation–Manifold Toolkit.
The purpose is to express:
\begin{itemize}
    \item reversible geometry via holonomy groups,
    \item curvature envelopes via sectional curvature,
    \item diagnostics via geometric projections,
    \item drift semigroup via geodesic flow,
    \item invariant extraction via holonomy-invariant projections,
    \item steering via parallel transport automorphisms,
    \item collision invariant via geometric fixed points,
    \item altitude transitions as holonomy-preserving maps.
\end{itemize}

\subsection*{H.1 Manifold Structure}

Let $(M,g)$ be a smooth Riemannian manifold with metric $g$.

Let $\nabla$ denote the Levi-Civita connection.

Let $\mathrm{Hol}(p)$ denote the holonomy group at point $p \in M$.

\subsection*{H.2 Altitude Geometry}

Altitude $A_0$ (Reversible Geometry):
\[
G : M \to M,\quad G \text{ is a diffeomorphism},\quad G^{-1} \text{ exists}.
\]

Altitude $A_1$ (Conceptual Metric):
\[
M(x,y) = g_x(v,w) \quad \text{for } v,w \in T_x M.
\]

Altitude $\Phi_7$ (Curvature Envelope):
\[
E(x) = K_x(\sigma),
\]
where $K_x(\sigma)$ is sectional curvature of plane $\sigma \subset T_x M$.

Diagnostics $D_1, D_2, D_3$:
\[
C,S,\Sigma : M \to M
\]
are geometric projections or deformations.

Altitude $\Phi_8$ (Drift Semigroup):
\[
D_t(x) = \gamma_x(t),
\]
where $\gamma_x$ is a geodesic with initial condition $x$.

Altitude $\Phi_9$ (Invariant Extraction):
\[
I(x) = \pi_{\mathrm{Hol}(x)}(x),
\]
a projection onto holonomy-invariant quantities.

Altitude $S_1$ (Steering):
\[
\Theta(x) = P_{\gamma}(x),
\]
parallel transport along curve $\gamma$.

Altitude $I_1$ (Collision Invariant):
\[
\Omega(x) \in \{0,1\}
\]
is a fixed point of a holonomy-invariant projection.

\subsection*{H.3 Holonomy-Preserving Transitions}

Define altitude transition map:
\[
\phi_{a,b} : \mathcal{O}(a) \to \mathcal{O}(b)
\]
such that:
\[
\phi_{a,b}(T) \text{ preserves holonomy}.
\]

Allowed transitions satisfy:
\[
\mathrm{Hol}(x) = \mathrm{Hol}(\phi_{a,b}(x)).
\]

Forbidden transitions violate holonomy preservation.

\subsection*{H.4 Curvature Envelope Formalism}

Curvature envelope:
\[
E(x) = K_x(\sigma)
\]
is altitude-pure because:
\[
E : M \to \mathbb{R}
\]
does not modify $M$.

Purity:
\[
P(\Phi_7) \iff E \text{ is curvature-only}.
\]

\subsection*{H.5 Projection Geometry for Diagnostics}

Diagnostics are geometric projections:
\[
C,S,\Sigma : M \to M.
\]

Collapse:
\[
C(x) = \pi_{\mathcal{S}}(x)
\]
onto a submanifold $\mathcal{S}$.

Soft deformation:
\[
S(x) = x + \epsilon v_x
\]
for small $\epsilon$.

Safety:
\[
\Sigma(x) = 
\begin{cases}
1, & x \in \mathcal{S}_{\mathrm{safe}}, \\
0, & \text{otherwise}.
\end{cases}
\]

\subsection*{H.6 Drift Semigroup as Geodesic Flow}

Define geodesic flow:
\[
D_t(x) = \gamma_x(t),
\]
where:
\[
\nabla_{\dot{\gamma}} \dot{\gamma} = 0.
\]

Semigroup property:
\[
D_{t+s}(x) = D_t(D_s(x)).
\]

\subsection*{H.7 Holonomy-Invariant Projection}

Invariant extraction:
\[
I(x) = \pi_{\mathrm{Hol}(x)}(x)
\]
is a projection onto holonomy-invariant quantities.

Thus:
\[
I(I(x)) = I(x).
\]

\subsection*{H.8 Steering as Parallel Transport Automorphism}

Steering:
\[
\Theta(x) = P_{\gamma}(x)
\]
is parallel transport along curve $\gamma$.

Automorphism:
\[
\Theta^{-1} \circ \Theta = \mathrm{id}.
\]

\subsection*{H.9 Collision Invariant as Fixed Point}

Collision invariant:
\[
\Omega(x) = 
\begin{cases}
1, & x \text{ is holonomy-safe}, \\
0, & \text{otherwise}.
\end{cases}
\]

Fixed point:
\[
\Omega(\Omega(x)) = \Omega(x).
\]

\subsection*{H.10 Firewall as Holonomy-Safe Submanifold}

Define:
\[
M_{\mathrm{safe}} = \{ x \in M \mid \Omega(x)=1 \}.
\]

Firewall:
\[
\phi_{a,b}(x) \in M_{\mathrm{safe}}
\iff
T(a,b)=1.
\]

\subsection*{H.11 Final Safety Theorem}

\textbf{Theorem.}  
A state $x$ is altitude-safe if and only if:
\[
x \in M_{\mathrm{safe}}.
\]

Equivalently:
\[
\Omega(x)=1.
\]

This completes the holonomy and projection geometry formal appendix.
